\chapter{Analyse du besoin}
\label{chap:besoin}

\section{Objectifs de la mission}
% TODO : reformuler le thème du stage en objectifs concrets —
% sécuriser l'accès au serveur, industrialiser le déploiement,
% héberger plusieurs applications web de façon isolée et supervisée.
[À rédiger.]

\section{Contraintes}
% TODO : contraintes techniques (matériel disponible, réseau local,
% budget nul/quasi nul -> outils open-source et gratuits), contraintes
% de temps (16 semaines), contraintes de sécurité.
[À rédiger.]

\section{Choix technologiques}
% TODO : justifier les choix (Docker/Docker Compose, Traefik, Jenkins,
% Node.js/Express, etc.) — un tableau récapitulatif est utile ici.
\begin{table}[h]
  \centering
  \begin{tabularx}{\textwidth}{@{}lX@{}}
    \toprule
    \textbf{Besoin} & \textbf{Solution retenue} \\
    \midrule
    Virtualisation & Proxmox VE \\
    Isolation des applications & Docker / Docker Compose \\
    Reverse proxy \& HTTPS & Traefik + Let's Encrypt \\
    Nom de domaine dynamique & DuckDNS \\
    Intégration / déploiement continu & Jenkins \\
    Pare-feu & UFW \\
    Protection brute-force SSH & Fail2ban \\
    Alerting & Webhooks Discord \\
    \bottomrule
  \end{tabularx}
  \caption{Correspondance besoin / solution technique}
\end{table}

\clearpage
